\input{sol_pre.tex}
\input{q_2024-10-12_1.tex}
\begin{document}
\input{sol_head.tex}

\Savol{1}
$2321=11\cdot211$: \ $11+211=222$.
\Javob{C}

\Savol{2}
$2022^{2026}\equiv2^{2026}\pmod5$; $2^k$ ning qoldiqlari $2,4,3,1$ davriy, $2026=4\cdot506+2$ $\Rightarrow$ qoldiq $4$.\\
$7^k$ ning oxirgi raqami $7,9,3,1$; $2025=4\cdot506+1$ $\Rightarrow$ $7$. \ $4\cdot7=28$.
\Javob{B}

\Savol{3}
$3$ stul $=0{,}9$ stol $\Rightarrow$ $1$ stul $=0{,}3$ stol, \ $4$ stul $=1{,}2$ stol — $120\%$.
\Javob{A}

\Savol{4}
Unumlar (ar/soat): $S=\frac1{80}$, \ $B=\frac1{60}-\frac1{80}=\frac1{240}$, \ $A=\frac1{40}-\frac1{80}=\frac1{80}$.
\[\frac1{80}+\frac1{240}+\frac1{80}=\frac7{240}\ \Rightarrow\ t=\frac{240}7=34\frac27.\]
\Javob{D}

\Savol{5}
$x^{\frac35+\frac12}y^{-2-\frac73}=x^{\frac{11}{10}}y^{-\frac{13}3}$.
\Javob{B}

\Savol{6}
$\left(\sqrt3+1\right)^2=2\left(2+\sqrt3\right)$, \ $\left(\sqrt7-1\right)^2=8-2\sqrt7=-2\left(\sqrt7-4\right)$: \ $2-2=0$.
\Javob{A}

\Savol{7}
$18-12\sqrt2=\left(2\sqrt3-\sqrt6\right)^2$ $\Rightarrow$ surat $\sqrt6\left(\sqrt2-1\right)$; \ maxraj $\left|2-2\sqrt2\right|=2\left(\sqrt2-1\right)$. \ Natija $\dfrac{\sqrt6}2$.
\Javob{C}

\Savol{8}
\[\left(\frac ca-\frac ab\right)+\left(\frac ab-\frac bc\right)+\left(\frac bc-\frac ca\right)=0.\]
\Javob{C}

\Savol{9}
$9$ ta had: $S_9=\frac{30+86}2\cdot9=522$. \ Qöyilganlar: $522-116=406$.
\Javob{C}

\Savol{10}
$b_2^3=5832$ $\Rightarrow$ $b_2=18$; \ $b_1+b_3=60$, $b_1b_3=324$ $\Rightarrow$ $b_1=54$, $b_3=6$, $q=\frac13$.
\[S=\frac{54}{1-\frac13}=81.\]
\Javob{C}

\Savol{11}
$c=-1$: \ surat $a^2(b+1)-b^2(a+1)=ab(a-b)+(a-b)(a+b)=(a-b)(ab+a+b)=2(a-b)$;\\
maxraj $a(b+1)^2-b(a+1)^2=ab(b-a)+(a-b)=(a-b)(1-ab)=2(a-b)$. \ Natija $1$.
\Javob{A}

\Savol{12}
$x\ne1$, \ $25x-5>0$ $\Rightarrow$ $x>\frac15$: \ $\left(\frac15;1\right)\cup(1;\infty)$.
\Javob{C}

\Savol{13}
$\frac1{x_1}+\frac1{x_2}=6$, $\frac1{x_1x_2}=7$: \ $\frac1{x_1^2}+\frac1{x_2^2}=36-14=22$, \ köpaytma $49$.\\
$t^2-22t+49=0$: \ $1-22+49=28$.
\Javob{B}

\Savol{14}
$\sqrt[4]{(x-2)^4}=|x-2|$: \ $x\ne2$ va $(x-4)(x+3)<0$ $\Rightarrow$ $-3<x<4$.\\
Butun yeçimlar: $-2,-1,0,1,3$ — yiğindi $1$.
\Javob{D}

\Savol{15}
II çorak: $\cos\alpha<0$, $\sin\alpha>0$. \ $1+\tg^2\alpha=\frac{25}{16}$ $\Rightarrow$ $\cos\alpha=-\frac45$, \ $\sin\alpha=\tg\alpha\cos\alpha=\frac35$.
\Javob{B}

\Savol{16}
$u=\log_4x$: \ $\log_x\sqrt{4x}=\frac12\left(\frac1u+1\right)$, \ $-\frac1u\ge0$ $\Rightarrow$ $u<0$.
\[\frac12\left(\frac1u+1\right)=\frac1{u^2}\ \Rightarrow\ u^2+u-2=0\ \Rightarrow\ u=-2,\quad x=\frac1{16}.\]
\Javob{C}

\Savol{17}
$t=2^x>0$: \ $2t(t+1)+15=11(t+1)$ $\Rightarrow$ $2t^2-9t+4=0$, \ $t=4$ yoki $t=\frac12$ $\Rightarrow$ $x=2$, $x=-1$. \ Yiğindi $1$.
\Javob{A}

\Savol{18}
$u=\sqrt[3]{x+4}$: \ $u^2+3u+4=2u+6$ $\Rightarrow$ $u^2+u-2=0$, \ $u=1$ yoki $u=-2$ $\Rightarrow$ $x=-3$, $x=-12$. \ Yiğindi $-15$.
\Javob{B}

\Savol{19}
$-2\le x\le5$; \ $x>0$, $x\ne1$, $x\ne2$; \ $\sin x\ne\frac\pi3$ har doim. \ Butun sonlar: $3,4,5$ — $3$ ta.
\Javob{C}

\Savol{20}
Çetdan teng uzoqlikdagi juftlar: $\sin k\alpha+\sin(16-k)\alpha=2\sin8\alpha\cos(8-k)\alpha$, \ $\cos k\alpha+\cos(16-k)\alpha=2\cos8\alpha\cos(8-k)\alpha$; örtadagi had $\sin8\alpha$ va $\cos8\alpha$.\\
Surat ham, maxraj ham bir xil yiğindiga köpaytirilgan $\sin8\alpha$ va $\cos8\alpha$: \ natija $\tg8\alpha$.
\Javob{D}

\Savol{21}
$t=\ln x$: \ $\displaystyle\int\ln t\,dt=t\ln t-t=\ln x\left(\ln\ln x-1\right)+C$.
\Javob{D}

\Savol{22}
$y_1$, $y_2$, $y_4$ — juft; \ $y_3$: $\ctg^3x$ toq, $\sqrt[3]{x^2}$ juft — juft emas. \ $3$ ta.
\Javob{C}

\Savol{23}
$12$ ta teng yonli uçburçak, uçidagi burçagi $30^\circ$:
\[S=12\cdot\frac12\cdot4^2\sin30^\circ=48.\]
\Javob{B}

\Savol{24}
$B(5;0)$, $D(0;4)$, $\overrightarrow{BE}=2\overrightarrow{BC}$ $\Rightarrow$ $E(5;8)$. \ $ABED$ — trapetsiya: $\frac{4+8}2\cdot5=30$.
\Javob{C}

\Savol{25}
$f'(b)=0$ ($b$ — maksimum nuqtasi); \ $c$ va $d$ da $f$ kamayadi: $f'(c)<0$, $f'(d)<0$; \ $g$ kamayuvçi va $g(0)=0$: $g(a)>0$, $g(c)>0$, $g(e)<0$.\\
a) $0\cdot g(e)=0$ — notöğri; \ b) $(-)\cdot(+)<0$ — notöğri; \ c) $(-)\cdot(+)<0$ — notöğri; \ d) $(-)\cdot(-)>0$ — töğri. \ $1$ ta.
\Javob{A}

\Savol{26}
\fig{\begin{tikzpicture}[scale=0.17,font=\small]
\coordinate (A) at (-12,0); \coordinate (C) at (12,0); \coordinate (B) at (0,16); \coordinate (D) at (0,9); \coordinate (M) at (0,0);
\draw (A)--(B)--(C)--cycle; \draw (A)--(D)--(C); \draw (B)--(D); \draw[axs] (D)--(M);
\node[below left] at (A) {$A$}; \node[below right] at (C) {$C$}; \node[above] at (B) {$B$}; \node[right] at (D) {$D$}; \node[aux,below] at (M) {$M$};
\node[aux,below] at (-6,0) {$12$};
\end{tikzpicture}}
$B$ va $D$ — $AC$ ning örta perpendikulyarida; $M$ — $AC$ örtasi, $AM=12$.
\[BM=\sqrt{400-144}=16,\qquad DM=\sqrt{225-144}=9,\qquad x=16-9=7.\]
\Javob{B}

\Savol{27}
\fig{\begin{tikzpicture}[scale=0.15,font=\small]
\coordinate (E) at (0,0); \coordinate (M) at (5,0); \coordinate (A) at (15,0); \coordinate (D) at (35,0); \coordinate (F) at (0,5); \coordinate (L) at (5,5); \coordinate (G) at (0,10); \coordinate (N) at (5,10); \coordinate (K) at (15,10); \coordinate (B) at (15,20); \coordinate (C) at (35,20); \coordinate (P) at (11.667,10);
\fill[okgreen!30] (F)--(N)--(C)--cycle;
\draw (E)--(D)--(C)--(B)--(A); \draw (E)--(G)--(N)--(K); \draw (F)--(L); \draw (M)--(N); \draw (F)--(N)--(C)--cycle; \draw[axs] (N)--(P);
\node[below] at (E) {$E$}; \node[below] at (M) {$M$}; \node[below] at (A) {$A$}; \node[below] at (D) {$D$}; \node[left] at (F) {$F$}; \node[left] at (G) {$G$}; \node[above] at (N) {$N$}; \node[above] at (B) {$B$}; \node[above right] at (C) {$C$}; \node[aux,below] at (P) {$P$}; \node[right] at (K) {$K$};
\end{tikzpicture}}
$EF=FL=LM=5$, \ $FG=GN=5$, \ $MN=AM=AK=10$, \ $KB=AK$ $\Rightarrow$ $AB=AD=20$.\\
$F$ dan $C$ gaça: gorizontal $35$, vertikal $15$. $N$ orqali $ED$ ga parallel çiziq $FC$ ni $P$ da kesadi ($P$ ham $F$ dan $5$ yuqorida): gorizontal masofa $\frac5{15}\cdot35=\frac{35}3$, \ $NP=\frac{35}3-5=\frac{20}3$.
\[S_{FNC}=S_{FNP}+S_{PNC}=\frac12\cdot NP\cdot15=50.\]
\Javob{A}

\Savol{28}
\fig{\begin{tikzpicture}[scale=0.32,font=\small]
\coordinate (B) at (-4,0); \coordinate (C) at (4,0); \coordinate (A) at (0,9.165); \coordinate (K) at (-2.4,3.666); \coordinate (M) at (0,0);
\draw (A)--(B)--(C)--cycle; \draw (0,2.619) circle (2.619); \fill (K) circle (3pt); \fill (M) circle (3pt); \draw[ax] (A)--(M);
\node[above] at (A) {$A$}; \node[below left] at (B) {$B$}; \node[below right] at (C) {$C$}; \node[below] at (M) {$M$};
\node[left] at (-1.3,6.6) {$6$}; \node[left] at (-3.3,1.8) {$4$};
\end{tikzpicture}}
$BM=BK=4$ ($BC$ kiçik tomon), $AB=10$, $BC=8$. \ $AM=\sqrt{100-16}=2\sqrt{21}$: \ $S=\frac12\cdot8\cdot2\sqrt{21}=8\sqrt{21}$.
\Javob{D}

\Savol{29}
\fig{\begin{tikzpicture}[scale=0.95,font=\small]
\draw (0,0) circle (1.8); \draw (0,0.9) circle (0.9); \draw (0,-0.9) circle (0.9); \draw (1.2,0) circle (0.6); \draw (-1.2,0) circle (0.6);
\foreach \p in {(0,0),(0,0.9),(0,-0.9),(1.2,0),(-1.2,0)} \fill \p circle (1.3pt);
\draw[axs] (0,0.9)--(1.2,0); \draw[ax] (0,0)--(1.2,0); \draw[ax] (1.2,0)--(1.8,0);
\node[below left] at (0,0) {$O$}; \node[left] at (0,0.9) {$B$}; \node[left] at (0,-0.9) {$D$}; \node[below] at (1.2,0) {$C$}; \node[above] at (-1.2,0) {$A$};
\node[aux,above right] at (0.6,0.45) {$R+r$};
\end{tikzpicture}}
Katta aylana radiusi $2R$ ($B$ va $D$ aylanalar uni içkaridan urinadi). $OC=2R-r$, $OB=R$, $BC=R+r$, $\angle BOC=90^\circ$:
\[(R+r)^2=R^2+(2R-r)^2\ \Rightarrow\ 2Rr=4R^2-4Rr\ \Rightarrow\ 6r=4R,\quad \frac rR=\frac23.\]
\Javob{D}

\Savol{30}
$E$ dan tekislikka $EF$ perpendikulyar tuşiramiz: $AC\parallel EF\parallel BD$, $ACDB$ — trapetsiya, $F\in CD$ va $DF:FC=BE:EA=1:2$:
\[d_E=\frac{2\cdot BD+1\cdot AC}3=\frac{12+12}3=8.\]
\Javob{A}

\Savol{31}
$A\cap B=\{3;7\}\subset C$ $\Rightarrow$ $(A\cap B)\cup C=C$ — $6$ ta element, $2^6=64$.
\Javob{A}

\Savol{32}
Raqamlar $\{6,7,8,9\}$; qoldiqlar: $6,9\to0$; $7\to1$; $8\to2$.\\
$(0,0,0)$: $8$; \ $(0,1,2)$: $3!\cdot2=12$; \ $777$, $888$: $2$. \ Jami $22$.
\Javob{C}

\newpage
\Savol{33}
Akvarium asosining tomoni $a$, piramida asosining tomoni $\frac a4$, yuzi $\frac{a^2}{16}$. Piramida hajmi suv sathining kötarilişiga teng hajm siqib çiqaradi:
\[\Delta h\cdot a^2=\frac13\cdot\frac{a^2}{16}\cdot h_p\ \Rightarrow\ \Delta h=\frac{h_p}{48}.\]

\textbf{33.} $\Delta h=\dfrac{384}{48}=8$ sm — \textbf{A}.

\textbf{34.} Suv sathi kamida $48$ sm bölişi kerak: $a=16$, piramida asosi $4\times4$:
\[V_{\text{suv}}=16^2\cdot48-\frac13\cdot16\cdot48=12288-256=12032\ \text{sm}^3\approx12\ \text{litr} — \textbf{D}.\]

\textbf{35.} $h_p=48\cdot0{,}5=24$ sm — \textbf{F}.
\Javob{33 — A, \ 34 — D, \ 35 — F}

\newpage
\Savol{36}
$t=|x|\ge0$: \ $t^3-2t^2-4t+3=0$; \ $t=3$ — ildiz: \ $(t-3)\left(t^2+t-1\right)=0$.
\[t=3,\qquad t=\frac{\sqrt5-1}2\quad\left(t=\frac{-1-\sqrt5}2<0\right).\]
$x=\pm3$, \ $x=\pm\dfrac{\sqrt5-1}2$.

a) Manfiy ildizlar $2$ ta.

b) Ildizlar juft-jufti bilan qarama-qarşi: yiğindi $0$.
\Javob{a) 2; \quad b) 0}

\newpage
\Savol{37}
$\cos x\left(\cos4x-1\right)+\sin^22x=0$, \ $\cos4x-1=-2\sin^22x$:
\[\sin^22x\left(1-2\cos x\right)=0.\]
$\sin2x=0$ $\Rightarrow$ $x=\frac{\pi k}2$; \ $\cos x=\frac12$ $\Rightarrow$ $x=\pm\frac\pi3+2\pi k$.

a) $[0;2\pi]$: $0,\ \frac\pi2,\ \pi,\ \frac{3\pi}2,\ 2\pi$ va $\frac\pi3,\ \frac{5\pi}3$ — $7$ ta.

b) $[0;\pi]$: $0+\frac\pi2+\pi+\frac\pi3=\dfrac{11\pi}6$.
\Javob{a) 7; \quad b) $\dfrac{11\pi}6$}

\newpage
\Savol{38}
$g(x)=(x+1)^2-7(x+1)=x^2-5x-6$. \ $f(x)=(x+1)\,g(x+2)=(x+1)\left(x^2-x-12\right)$.

a) $g(x)=\left(x-\frac52\right)^2-\frac{49}4\ge-12{,}25$. \ Eng kiçik butun qiymat $-12$.

b) $f-g=(x+1)\left(x^2-x-12-x+6\right)=(x+1)\left(x^2-2x-6\right)=x^3-x^2-8x-6$.
\[x^3-x^2-8x+8=0\ \Rightarrow\ (x-1)\left(x^2-8\right)=0:\quad 1+2\sqrt2-2\sqrt2=1.\]
\Javob{a) $-12$; \quad b) 1}

\newpage
\Savol{39}
Urinma $Ox$ bilan $135^\circ$ hosil qiladi: $k=f'(4)=\tg135^\circ=-1$, \ $f(4)=3$.

a) $g(2x+1)=\left(x^2+x+1\right)f(x)$ ni differensiallaymiz:
\[2g'(2x+1)=(2x+1)f(x)+\left(x^2+x+1\right)f'(x).\]
$x=4$: \ $2g'(9)=9\cdot3+21\cdot(-1)=6$ $\Rightarrow$ $g'(9)=3$.

b) $y=-x+b$, $(4;3)$ dan ötadi: $b=7$. \ $k+b=6$.
\Javob{a) 3; \quad b) 6}

\newpage
\Savol{40}
\fig{\begin{tikzpicture}[x=0.6cm,y=0.45cm,font=\small]
\fill[okgreen!30] plot[domain=1:5,samples=40] (\x,{\x*\x-6*\x+5})--(3,-4)--cycle;
\draw[-{Stealth[length=4pt]}] (-0.5,0)--(6.8,0) node[right] {$x$}; \draw[-{Stealth[length=4pt]}] (0,-4.8)--(0,5.8) node[left] {$y$};
\draw[thick,domain=-0.2:6.2,samples=60] plot (\x,{\x*\x-6*\x+5}); \draw[thick] (-0.3,2.6)--(3,-4)--(6.3,2.6);
\node[below] at (1,0) {$1$}; \node[below] at (5,0) {$5$}; \node[left] at (3,-4.3) {$(3;-4)$};
\end{tikzpicture}}
$u=x-3$: \ $f=u^2-4$, \ $g=2|u|-4$.

a) $u^2=2|u|$ $\Rightarrow$ $u=0,\ \pm2$ $\Rightarrow$ $x=1,\ 3,\ 5$ — $3$ ta.

b) $|u|\le2$ da $g-f=2|u|-u^2\ge0$; simmetriya:
\[S=2\int_0^2\left(2u-u^2\right)du=2\left(4-\frac83\right)=\frac83.\]
\Javob{a) 3; \quad b) $\dfrac83$}

\newpage
\Savol{41}
\fig{\begin{tikzpicture}[scale=0.65,font=\small]
\coordinate (C) at (0,2); \coordinate (B) at (6,0); \coordinate (A) at (1,5); \coordinate (P) at (4,2);
\draw (0,0) rectangle (2,2); \draw (2,0) rectangle (4,2); \draw (4,0) rectangle (6,2);
\draw (A)--(B)--(C)--cycle; \draw[axs] (C)--(P); \fill (P) circle (1.5pt);
\draw ($(C)!0.3cm!(B)$)--($(C)!0.3cm!(B)!0.3cm!90:(B)$)--($(C)!0.3cm!(A)$);
\node[above] at (A) {$A$}; \node[right] at (B) {$B$}; \node[left] at (C) {$C$}; \node[aux,above right] at (P) {$P$};
\end{tikzpicture}}
$AB$ örta kvadratning $P$ uçidan ötadi; $PB$ — kvadrat diagonali, gorizontal bilan $45^\circ$.\\
$CB$: gorizontal $6$, vertikal $2$ $\Rightarrow$ $CB=\sqrt{40}=2\sqrt{10}$, gorizontal bilan burçagi $\theta$, $\tg\theta=\frac13$.
\[\tg\angle ABC=\tg\left(45^\circ-\theta\right)=\frac{1-\frac13}{1+\frac13}=\frac12\ \Rightarrow\ AC=CB\cdot\frac12=\sqrt{10}.\]

a) $AB=\sqrt{40+10}=5\sqrt2$.

b) $S=\frac12\cdot\sqrt{10}\cdot2\sqrt{10}=10$.
\Javob{a) $5\sqrt2$; \quad b) 10}

\newpage
\Savol{42}
\fig{\begin{tikzpicture}[scale=1.5,font=\small]
\coordinate (O) at (0,0); \coordinate (A) at (0,1); \coordinate (B) at (0.866,1.5); \coordinate (C) at (1.732,1); \coordinate (D) at (1.732,0); \coordinate (E) at (0.866,-0.5);
\fill[okgreen!30] (A)--(C)--(E) arc (-30:90:1)--cycle;
\draw (O) circle (1); \draw (O)--(A)--(B)--(C)--(D)--(E)--cycle; \draw (A)--(C)--(E); \draw[ax] (A)--(E);
\fill (O) circle (1pt); \node[left] at (O) {$O$}; \node[above left] at (A) {$A$}; \node[above] at (B) {$B$}; \node[right] at (C) {$C$}; \node[right] at (D) {$D$}; \node[below] at (E) {$E$};
\end{tikzpicture}}
Oltiburçak tomoni $a=R$ ($O$ — oltiburçak uçi, $OA=OE=R$), $\angle AOE=120^\circ$; sektor $AOE$ töliq oltiburçak içida.\\
$a^2=R^2=\sqrt3+1$, \ $\pi=3$.

a) Içkarida: $S_{\text{sek}}=\frac13\pi R^2=R^2$. \ Oltiburçak: $\frac{3\sqrt3}2R^2$. \ Taşqarida: $\left(\frac{3\sqrt3}2-1\right)R^2$:
\[\frac{3\sqrt3-2}2.\]

b) $\triangle ACE$ — tomoni $R\sqrt3$ bölgan teng tomonli: $S=\frac{3\sqrt3}4R^2$. Undan $AE$ vatar ajratgan segment ayiriladi:
\[S_{\text{seg}}=\frac{\pi R^2}3-\frac12R^2\sin120^\circ=R^2-\frac{\sqrt3}4R^2.\]
\[S=\frac{3\sqrt3}4R^2-R^2+\frac{\sqrt3}4R^2=\left(\sqrt3-1\right)R^2=\left(\sqrt3-1\right)\left(\sqrt3+1\right)=2.\]
\Javob{a) $\dfrac{3\sqrt3-2}2$; \quad b) 2}

\newpage
\Savol{43}
\fig{\begin{tikzpicture}[scale=0.17,font=\small]
\coordinate (D) at (0,0); \coordinate (B) at (-5.352,20.49); \coordinate (A) at (-10.247,0); \coordinate (C) at (53.43,0);
\coordinate (K) at (-5.179,0); \coordinate (L) at (6.179,0); \coordinate (O) at (-5.179,4); \coordinate (N) at (6.179,8); \coordinate (J) at (-1.561,5.978); \coordinate (I) at (-1.309,5.011);
\draw (A)--(B)--(C)--cycle; \draw (B)--(D); \draw (O) circle (4); \draw (N) circle (8);
\draw[axs] (O)--(D)--(N); \draw[ax] (O)--(K); \draw[ax] (N)--(L);
\foreach \p in {D,K,L,O,N,I,J} \fill (\p) circle (7pt);
\node[above] at (B) {$B$}; \node[below] at (A) {$A$}; \node[below] at (C) {$C$}; \node[below] at (D) {$D$}; \node[below] at (K) {$K$}; \node[below] at (L) {$L$}; \node[left] at (O) {$O$}; \node[above] at (N) {$N$}; \node[left] at (-1.4,4.4) {$I$}; \node[left] at (-2.0,7.3) {$J$};
\end{tikzpicture}}
Urinma kesmalari: $DK=DI$, $DL=DJ$, $BI=BJ+JI$ (yoki aksinça).\\
$DO$ va $DN$ — qöşni burçaklar bissektrisalari: $\angle ODN=90^\circ$ $\Rightarrow$ $\angle ODK=\angle DNL$, \ $\triangle ODK\sim\triangle DNL$:
\[\frac{DK}{4}=\frac8{DL}\ \Rightarrow\ DK\cdot DL=32,\qquad DK+DL=KL=\sqrt{129}.\]

a) $IJ=|DL-DK|=\sqrt{(DK+DL)^2-4DK\cdot DL}=\sqrt{129-128}=1$.

b) $\angle ABD=x$, $\angle DBC=2x$, $t=\tg\frac x2$. \ $BI=4\ctg\frac x2=\frac4t$, \ $BJ=8\ctg x=\frac{4\left(1-t^2\right)}t$.\\
$BI-BJ=\frac{4t^2}t=4t>0$ va $BI-BJ=IJ=1$ $\Rightarrow$ $t=\frac14$:
\[BJ=\frac{4\left(1-\frac1{16}\right)}{\frac14}=15.\]
\Javob{a) 1; \quad b) 15}

\newpage
\Savol{44}
\fig{\begin{tikzpicture}[scale=1.1,font=\small]
\draw (-1.5,0) arc (180:360:1.5); \draw (-1.5,0)--(1.5,0);
\fill[okgreen!30] (-1.06,0) rectangle (1.06,-1.06); \draw (-1.06,0) rectangle (1.06,-1.06);
\draw[axs] (0,0)--(1.06,-1.06); \fill (0,0) circle (1pt);
\node[above] at (0,0) {$O$}; \node[aux,right] at (0.6,-0.5) {$r$}; \node[above] at (0.53,0) {\scriptsize$x$}; \node[right] at (1.06,-0.53) {\scriptsize$y$};
\end{tikzpicture}}
Kesim: radiusi $r$ bölgan yarim doiraga içki çizilgan töğri törtburçak $2x\times y$, $x^2+y^2=r^2$. Uzunlik umumiy, demak $2xy$ maksimal bölişi kerak:
\[2xy\le x^2+y^2=r^2,\quad \text{tenglik } x=y=\frac r{\sqrt2}\ \text{da}.\]
Parallelepiped asosi (kesimi): $r\sqrt2\times\frac r{\sqrt2}$, diagonali $\sqrt{2r^2+\frac{r^2}2}=\dfrac{r\sqrt{10}}2$.

a) $\dfrac{2r}{r\sqrt{10}/2}=\dfrac4{\sqrt{10}}=\dfrac{2\sqrt{10}}5$.

b) $\dfrac{r^2\cdot l}{\frac12\pi r^2\cdot l}=\dfrac2\pi=\dfrac23$.
\Javob{a) $\dfrac{2\sqrt{10}}5$; \quad b) $\dfrac23$}

\newpage
\Savol{45}
\fig{\begin{tikzpicture}[scale=0.8,font=\small]
\draw (-0.5,0)--(5.5,0);
\begin{scope}[shift={(1,0)},rotate=30]
\coordinate (P) at (0,0); \coordinate (Q) at (4,0); \coordinate (S) at (0,1.6); \coordinate (T) at (4,1.6);
\fill[okgreen!30] (P)--(Q)--(1.229,1.6)--(S)--cycle;
\draw (P)--(Q)--(T)--(S)--cycle; \draw[axs] (Q)--(1.229,1.6);
\node[aux,above] at (2.6,1.6) {$2\sqrt3R$}; \node[left] at (0,0.8) {$2R$};
\end{scope}
\draw (1,0)+(0:0.7) arc (0:30:0.7); \node at (2.05,0.2) {\scriptsize$30^\circ$};
\end{tikzpicture}}
Öq gorizontal bilan $30^\circ$ hosil qiladi; suv sirti gorizontal, demak u silindr asosi bilan $60^\circ$ burçak hosil qiladi. Suv yuqori asosning eng past nuqtasigaça qoladi; qarama-qarşi yasovçida sirt $2R\tg60^\circ=2\sqrt3R$ pastda.\\
Tökilgan qism — silindrning qiya kesilgan bölagi, hajmi örtaça balandligi $\sqrt3R$ bölgan silindrga teng:
\[h=H-\sqrt3R,\qquad \frac hH=1-\frac{\sqrt3R}H=1-\sqrt3\cdot\frac{\sqrt3}8=\frac58\quad\left(\frac RH=\frac{\sqrt3}8\right).\]

a) $\dfrac hH=\dfrac58$.

b) Tökilgani: $\dfrac38\cdot400=150$ litr.
\Javob{a) $\dfrac58$; \quad b) 150}

\ShowJavoblar
\end{document}
